\subsection{From the two-square marginal to the square}
\label{app:pair}
To reintroduce a diagonal restriction, we need the probability that the two rooks occupy its
endpoints. We estimate it first with the ratio of predictors $\pi^T$ and then with the local
expression $\pi_1$ on the starting board. Let $\Pi$ be the orthogonal projector onto the
row-minus-column mode. At the rook stationary point, set
$G^0=(S_T+\Pi)^{-1}$, $\mathsf G(i,j)=a_i^\top G^0a_j$. The stationary class requires
$\max p_e\le C/k$ and $\lambda_{\min}(S_T+\Pi)\ge\gamma_A>0$. The scaling class above implies
these margins for large $k$. For a pair $E=\{e,f\}$ in distinct rows and columns, define
\[
 \pi^T=w_ew_f e^{L_T(M_{ef})-L_T(M)},\quad
 \pi_1=p_ep_f\{1-\mathsf G(e,f)+\hat c_e+\hat c_f\},\quad
 X=p_ep_f(p_e+p_f)\mathsf G(e,f).
\]
On the binary mask, the ratio $\pi^{\rm ex}/\pi^T$ is exactly
$\exp\{g_T(M_{ef})-g_T(M)\}$. Removing two fibers from each family changes their sums by
$O(k^{-1})$. The dual system in the fixed gauge restores the masses with potentials of order
$O(k^{-1})$. The gap decreases by at most $O(k^{-1})$. By the Schur complement and the
resolvent identity, the retained entries of the Green matrices change by $O(k^{-2})$.
Omitting the indices in the normalized contractions gives
$|\mathfrak F(M_{ef})-\mathfrak F(M)|=O(k^{-1})$, where $\mathfrak F$ is evaluated at the
profile of each mask. Consequently $\Delta_E g_T=g_T(M_{ef})-g_T(M)=O(k^{-2})$ by (T).
Two applications of rook insertion and the stability of $p,\hat c$ give $\pi^T=\pi_1+O(k^{-4})$
(Corollary C for rooks). This gives $\pi^{\rm ex}-\pi_1=O(k^{-4})$, which is statement (R0).

\paragraph{The third square [D].}
When we occupy a third square $c$ in another row and another column, the pair remains available
in the binary minor $M_c$. Applying (R0) to both masks would give only $O(k^{-4})$ for the
difference. Compare $M,M_{ef},M_c,(M_c)_{ef}$ in $\Delta_c\Delta_E(\Phi/k)$ using two mass
parameters. Removing $E$ has parameter $s$, removing $c$ has parameter $t$, and the inverse
size is $k^{-1}/(1-2s-t)$. The kernels and inverses defining $\Phi$ are bounded on fixed
complex disks. The expression divided by $1/k$ has a bounded mixed derivative there.
Integrating over $[0,1/k]^2$ gives $|\Delta_c\Delta_E(\Phi/k)|=O(k^{-3})$. The same bound for
the remainder follows from two applications of $(T^+)$. Then, in
\[
 \pi^{\rm ex}-\pi_1=(\pi^T-\pi_1)+\pi^T(\exp(\Delta_E g_T)-1)
\]
the second summand changes by $O(p_ep_f/k^3)$ by the product rule and
$\Delta_c\pi^T=O(p_ep_f/k)$. For the first we retain $p_ep_f$ in the product of the insertions.
The second differences of $\mathsf G,\hat c$ and the change of order $O(k^{-3})$ in $r_e$
in a distant minor give $|\Delta_c(\pi^T-\pi_1)|\le Cp_ep_f/k^3$. This proves (R1) with weight
$\varpi(E)=k^2p_ep_f$. We now average over the square occupied in a row $r$ distinct from the
rows of the pair. Commutation of three minors and $\sum_{c\in r}\pi^{\rm der}_c=1$ give the
consistency defect of $\pi_1$ under this average,
\[
  D_r=\sum_{\substack{c\in r\\ M_c\supset E}}\pi^{\rm der}_c(M)\pi_1(M_c)-\pi_1(M),
       \qquad |D_r|\le C\varpi(E)/k^5.
\]
To prove this, set $R_\pi=\pi^T-\pi_1$. Commutation gives three summands in $D_r$.
The first is $\pi^T(M)\sum_c\pi_c^{\rm der}(M_{ef})(e^{\Delta_E r_c}-1)$.
Stability (S) of the insertion error gives $\Delta_E r_c=O(k^{-3})$.
The second is $-\sum_c\pi_c^{\rm der}(M)\Delta_cR_\pi$ over the squares that retain the pair.
The preceding increment bound applies here. The third,
$R_\pi(M)\sum_{c\text{ lost}}\pi_c^{\rm der}(M)$, involves at most two squares, each of
mass $O(1/k)$. Combining the three bounds gives (R2).

On binary masks the gap and the factorization of the stationary point give a lower bound
depending only on the class for $kp_e$ on the support, so that absolute bounds for a fixed
number of pairs can be converted into weighted bounds. Extending this conversion to arbitrary
forbidden squares or repeating it for growing groups would require another estimate.
Two collections of blocks with a third vertex have $\hat c_e=\hat c_f=\mathsf G(e,f)=0$ and
$\pi^{\rm ex}=\pi_1$ exactly [O]. In these families the error vanishes for any positive value
of the constant in the bound.

\paragraph{A pair and a square [D].}
On the board $B$ that restores the diagonal restriction of a single line $E$, the logarithm of
the count changes by $\log(1-\pi^{\rm ex})$. To subtract the change in the predictor from this,
let $\lambda=L(B)-L_T(M)$. With $p=p_e,q=p_f$, the first expansion gives
$\lambda=-\pi_1+O(k^{-4})$ (Theorem 1). The second (Theorem 2) gives
\[
       \lambda=\log(1-\pi_1)+X+O(k^{-5}).
\]
At the trial point, a rank-two update of the axial operator incorporates the two squares, and
a Schur complement incorporates the line. The Newton displacement has component $pq+O(k^{-3})$
on the four fibers and $2pq+O(k^{-3})$ on the added line.
The first variation and half the second variation of the action, the local factor, and the
trace of the logarithm leave an analytic function of seven variables of order $O(k^{-1})$.
The degree-four polynomial for $\lambda+\pi_1$ is
$-p^2q^2/2+pq(p+q)\mathsf G(e,f)$ [O, rational coefficients].
The denominators at the origin equal one and the Taylor bound is uniform. Thus the error
$g^\sharp=\log(1-\pi^{\rm ex})-\lambda$ retains the difference between the two logarithms and
$-X$ (Corollary 5).

The average over pairs already includes the effect of each diagonal separately. If all four
squares are available, it remains to measure what imposing both restrictions at once adds. Let
$E'$ be the other pair and $F_\square=E\cup E'$. The update on four squares and the two-dimensional
Schur complement~\cite{hj} give a function of eighteen small arguments.
After subtracting the two single-line expressions taken on the same mask, its degree-four
polynomial is $2\prod_{i\in F_\square}p_i$ [O]. This is the predictor excess (Theorem 3), with
remainder $O(k^{-5})$. The two pairs are incompatible and leave the same minor. If the completion
counts are positive, the exact avoidance excess is $\log[(1-2\pi^{\rm ex})/(1-\pi^{\rm ex})^2]$.
Its first term is $-(\pi^{\rm ex})^2=-\prod_{i\in F_\square}p_i+O(k^{-5})$ by (R0) and by the
factorization of $p_i(1-p_i)$. The joint contribution to the error is therefore
$-3\prod_{i\in F_\square}p_i+O(k^{-5})$ (Corollary 6). Taking the boundary mean counts this
difference once per complete square and gives $-S$ in~\eqref{eq:main}.
